\section{Elementary analytic primitives with finite bounds}
\label{app:primitives}

This appendix supplies explicit definitions and sufficient bounds for the
analytic notation in the report. They are deliberately generous. Their
role is to make primitive recursiveness and the finite-error proof
translation transparent, rather than to optimize numerical evaluation.

\subsection{Sine, cosine, exponential, and the constants}

Use the usual power-series definitions, understood as rational
approximation procedures with the following remainders. For rational $z$,
let $M=\max(1,\lceil|z|\rceil)$ and
\[
 S_N(z)=\sum_{j=0}^{N-1}\frac{(-1)^jz^{2j+1}}{(2j+1)!}.
\]
Set $N=2M^2+k+2$. The omitted alternating tail starts in a decreasing
region, and its absolute value is at most
\[
 \frac{M^{2N+1}}{(2N+1)!}
 \le\left(\frac{M^2}{N+1}\right)^{N+1}
 \le2^{-(N+1)}\le2^{-k}.
\]
For the middle inequality, the final $N+1$ factors of $(2N+1)!$ are
at least $N+1$, and $M^{2N+1}\le M^{2N+2}$. Consecutive omitted terms
decrease since their ratio is at most
$M^2/((2N+2)(2N+3))<1$. Negative arguments follow by oddness.
Thus this is a single explicitly bounded rational computation.

For $\exp z$ with $|z|\le M$, put $m=2M^2+k+3$ and sum its series
through degree $2m-1$. The first omitted absolute term satisfies
\[
 \frac{M^{2m}}{(2m)!}
 \le\left(\frac{M^2}{m}\right)^m\le2^{-m}.
\]
Subsequent ratios are at most $1/2$, so the whole omitted absolute tail
is at most $2^{1-m}\le2^{-k}$. The same absolute-series method gives
cosine approximations. Factorials, finite powers, and rational sums are
primitive recursive; no analytic convergence test is executed.

For $\log2$, use
\[
 \log2=\sum_{j=1}^{\infty}\frac1{j2^j}.
\]
The tail after $N$ terms is at most $2^{-N}$. This suffices for every
precision required by the full language, although the reduction itself
does not use this constant.

For $\pi$, use Machin's formula
\begin{equation}\label{eq:machin}
 \pi=16\arctan(1/5)-4\arctan(1/239),
 \qquad
 \arctan r=\sum_{j=0}^{\infty}\frac{(-1)^jr^{2j+1}}{2j+1}.
\end{equation}
Truncating both arctangent series after $N=k+2$ terms has combined error
at most
\[
 20\,5^{-2N-1}\le2^{-k}.
\]
The alternating lower and upper bounds already give
\[
 16\left(\frac15-\frac1{3\cdot5^3}\right)-\frac4{239}
 <\pi<\frac{16}{5},
\]
and the rational number on the left is greater than $3$.
These are finite rational certificates of the numerical bounds used in
the rounding argument.

\subsection{The elementary identities needed by the construction}

The sine/cosine addition formulas and
$\sin^2x+\cos^2x=1$ follow from multiplication of their power series.
For an arithmetical proof, multiply finite sums, prove the finite
coefficient identities, and bound all omitted terms by the absolute
series estimates above. An arbitrary requested final precision determines
the truncation indices. Thus the usual series proof has a direct finite
implementation; rearrangement of an unbounded sum is unnecessary.

One can verify the relevant consequences of Machin's formula without an
unbounded search for a trigonometric zero. Put
$a=\arctan(1/5)$ and $b=\arctan(1/239)$. The addition identities give
\[
 \tan(2a)=\frac5{12},\qquad
 \tan(4a)=\frac{120}{119},\qquad
 \tan(4a-b)=1.
\]
The last identity is the exact rational calculation
\[
 \frac{120/119-1/239}{1+(120/119)(1/239)}=1.
\]
For $\theta=4a-b$, the series bounds give $0<\theta<1$;
the sine and cosine of $\theta$ are positive. Thus
$\sin\theta=\cos\theta$. The circle and addition identities yield
$\sin(2\theta)=1$, $\cos(2\theta)=0$,
$\sin(4\theta)=0$, and $\cos(4\theta)=-1$.
Since \eqref{eq:machin} is $\pi=4\theta$, we obtain exactly
\[
 \sin\pi=0,\quad\cos\pi=-1,\quad
 \sin(\pi/2)=1,\quad\sin(3\pi/2)=-1.
\]
Integer-shift periodicity follows by finite induction using addition.

The relation $\tan(\arctan r)=r$ on the small rational interval in use
also has a finite-series proof. Formally differentiating the series for
$A(r)=\arctan r$ gives $A'(r)=1/(1+r^2)$. If
$B(r)=\sin A(r)-r\cos A(r)$, the addition/composition series imply
$B'(r)=rB(r)/(1+r^2)$ and $B(0)=0$. Coefficient induction gives zero
for every coefficient of $B$. On $|r|\le1/5$, geometric bounds for
the tails justify the resulting finite-precision identity. All
coefficient manipulations concern rational polynomials and bounded sums;
no general differential-equation existence or uniqueness theorem is used.

We also use
\begin{equation}\label{eq:sine-basic-bounds}
 |\sin x|\le1,\qquad
 |\sin x-\sin y|\le|x-y|,\qquad
 \sin s\ge s-s^3/6\quad(0\le s\le2).
\end{equation}
The first follows from the circle identity. The last follows from the
alternating series, whose terms decrease on $[0,2]$. For the Lipschitz
bound, the addition formulas give the finite-increment estimate
\[
 |\sin(x+h)-\sin x|\le |h|+h^2
 \qquad(|h|\le1),
\]
using $|\sin h|\le|h|$ and $|1-\cos h|\le h^2/2$ on that interval,
both immediate from their alternating series. Subdivide $[x,y]$ into
$r$ equal steps, with $r\ge|x-y|$. Telescoping gives
\[
 |\sin x-\sin y|\le |x-y|+\frac{|x-y|^2}{r}.
\]
Given a rational error budget $2^{-k}$, obtain an integer magnitude bound
$M\ge\max(1,|x-y|)$ from the available rational enclosures and choose
$r\ge\max(M,2^kM^2)$. This proves the Lipschitz inequality in
finite-precision form. The proof uses only explicitly bounded finite
telescoping, not a mean-value theorem with an unspecified real witness.

\subsection{A global bound for evaluating arbitrary expression syntax}

After composition has been expanded, let $N$ be the number of expression
nodes. Choose an integer $B_0\ge8$ dominating the absolute rational input
and all rational constants in the syntax. Define
\[
 B_{j+1}=2^{4B_j},\qquad B=B_N.
\]
Every exact subexpression value has magnitude at most $B$. Addition and
multiplication are covered by the recurrence, sine by $1$, absolute value
by its argument bound, and exponentiation by
$\exp(B_j)\le 2^{2B_j}<B_{j+1}$. The constants $\pi$ and $\log2$ are
covered by $B_0$. The elementary bound for the exponential follows from
$e<4$ and monotonicity, or directly from its series.

For child errors at most $1$, the local operations have a common error
amplification bound
\[
 L=2^{4(B+1)}.
\]
For addition and multiplication this follows from their elementary error
identities; for sine and absolute value it follows from their Lipschitz
bounds. For exponential on $[-B-1,B+1]$, its series or a finite-increment
estimate gives amplification at most $\exp(B+1)<L$.

Evaluate nodes from children to parents, using local approximations of
error at most $2^{-K}$, where
\begin{equation}\label{eq:global-precision}
 K=k+(N+1)\bigl(4(B+1)+2\bigr)+4.
\end{equation}
At depth $j$, the accumulated error is bounded by
\[
 2^{-K}(1+L+\cdots+L^j)\le2^{1-K}L^N\le2^{-k}.
\]
The same bound is below $1$, so the hypothesis used for each local
amplification estimate is maintained. A final outward rounding, if
desired, fits within the generous extra factors in
\eqref{eq:global-precision}. Alternatively leave the finite rational
value unrounded.

This constructs the function $\App$ of Lemma~\ref{lem:approx}. The
variable-height tower $B_N$ is primitive recursive because its height is
a finite recursion parameter for one fixed exponential operation. The
recursion does not evaluate arbitrary primitive recursive programs and
does not contradict the absence of a primitive recursive universal
interpreter for that larger class.

The much smaller reference implementation in the archive handles only the
language actually used by the reduction. It uses interval bounds adapted
to each expression node, instead of the very large uniform tower. Both
approaches have explicitly predetermined finite precision bounds.
